\documentclass[11pt]{article}
\usepackage[margin=1in]{geometry}
\usepackage{amsmath,amssymb}
\usepackage{booktabs}
\usepackage{array}
\usepackage{setspace}
\usepackage{enumitem}
\usepackage[authoryear,round]{natbib}
\usepackage[hidelinks]{hyperref}
\usepackage{xurl}
\setstretch{1.12}
\setlist{itemsep=2pt,topsep=4pt}

\title{Statistically Detectable, Economically Small:\\
A Pre-Registered Out-of-Sample Test of Minute-Level\\ Bitcoin Return Predictability}
\author{Gaoshan Zeng\thanks{Independent researcher. Email: zenggaoshan05@gmail.com. Code, frozen protocols, execution contracts and evidence files: \url{https://github.com/zenggs05/QuantForge} (source-available; noncommercial license). Market data: Binance Vision public data (\url{https://data.binance.vision}), used under its dataset terms. The use of AI tools is described in the disclosure statement at the end of the paper.}}
\date{Preprint v1, November 2026. Comments welcome.}

\begin{document}
\maketitle

\begin{abstract}
\noindent
We test whether a fixed forecasting program built from public one-minute candles of Binance BTCUSDT spot carries out-of-sample information about 15-minute returns, and how much that information is worth. The program (fourteen causal features, Ridge regression, expanding window) and its tests were frozen and hashed before the evaluation data were read. A historical external validation over 56 monthly folds (2021--2025, 2.45 million minutes) rejects the null of no predictive information ($t=5.44$), but the evidence is uneven: 76\% of the statistic comes from 2021 (excluding 2021, $t=3.00$), and a test on the 56 monthly means gives $t=2.14$. A within-month demeaned test ($t=5.74$) shows that the forecasts carry timing information rather than agreement with the sample trend. Two confirmatory holdouts, each run under a protocol frozen before its data were opened, split. July--August 2026 rejects ($t=1.89$, $p=0.029$). September--October 2026, with a drift-free primary statistic, does not ($t=1.33$, $p=0.092$), although its estimate (an information coefficient of 0.021) is close to the historical one (0.024). Economically the information is small: squared-error tests never reject, the 95\% upper bound on the out-of-sample $R^2$ is about 0.15--0.2\%, and a non-overlapping 15-minute rule breaks even at 0.2--0.3 basis points per side, against a retail taker fee of 10 basis points. A census of 2,070 days of Binance archives finds 157 minutes where the exchange's monthly and daily files disagree; excluding them leaves the results unchanged. Short-horizon predictability of BTCUSDT from public candles is weak, fragile in short samples and economically small.
\end{abstract}

\medskip
\noindent\textbf{Keywords:} return predictability; cryptocurrency; out-of-sample forecasting; Clark--West test; pre-registration; transaction costs; data quality.\\
\noindent\textbf{JEL codes:} C12, C53, C58, G14, G17.

\newpage
\section{Introduction}

High-frequency cryptocurrency data are cheap and abundant, and so is the scope for false discoveries. Researchers can choose among many features, horizons, sample splits, model classes and tuning parameters, and the same data can be revisited until something looks significant \citep{white2000,harvey2016,harvey2017}. This paper asks a deliberately narrow question. Does a fixed program, built only from public one-minute candles and fully specified in advance, carry information about the next 15 minutes of BTCUSDT returns on data it has never seen? And if so, how much is that information worth?

We answer with tests whose design was frozen before the corresponding data were touched. The features, label, timing conventions, model and hyperparameters of the forecasting program were fixed before any data were downloaded. Ten months of development data (September 2025--June 2026) were then used, in exploratory analyses, to choose the primary test and its bandwidth, to define a secondary specification and to plan power; the source files were then fingerprinted. A pre-registered historical extension downloaded 68 months of earlier data, audited them without computing any forecast, and evaluated the frozen program in 56 strictly forward-rolling monthly folds. A two-month holdout (July--August 2026) had been sealed by project policy from the start. Code guards enforced the seal from 4 November 2026, when the confirmatory protocol was frozen and hashed. The holdout was authorized twelve days later, in writing and citing that hash, and opened the following day. A second holdout (September--October 2026) was then tested under its own protocol, frozen before its data were accessed. Every model-based computation was executed twice in independent processes, and the deterministic outputs had to match byte for byte.

The results fall into separate evidence layers. On the \emph{predictive-information} layer, measured with the \citet{clarkwest2007} test for nested models, the historical test and the first holdout reject the null of no information. The historical rejection is strong but uneven over time; the first holdout's is narrow but becomes stronger once the forecast's mean is removed. The second holdout, tested with a drift-free information coefficient, does not reject ($p=0.092$), although its point estimate is close to the historical one. On the \emph{forecast-accuracy} layer, measured with mean squared error and the \citet{dieboldmariano1995} test, there is no reliable gain over a zero forecast, and confidence bounds rule out an out-of-sample $R^2$ above roughly 0.15--0.2\%. On the \emph{economic} layer, a gross-of-cost analysis of the stored forecasts shows that the edge per trade is a fraction of a basis point for unconditional trading. Historically and in Holdout A it reaches a few basis points only for forecasts above 5 basis points, about 1.3 percent of periods; in Holdout B the selective rules do not break even at all. Both figures are well below retail fees, and we do not model latency, queue position or adverse selection.

The paper makes three contributions. First, it gives a clean out-of-sample answer to a question that is usually answered with many researcher degrees of freedom: short-horizon predictability in BTCUSDT spot from public candles is statistically detectable but fragile in short samples, and economically small, consistent with the view that liquid markets can be predictable in a statistical sense while remaining efficient at the margin where predictability could be exploited \citep{chordia2008}. Second, it shows how to separate predictive information from forecast accuracy and from agreement with the sample drift, using drift-free variants of the test, time-shift placebos and calibration regressions. Third, it documents a research protocol suited to public high-frequency data: frozen plans, hashed execution contracts, result-blind data audits, byte-identical reproduction, and an explicit ledger of what was known before each test. It also reports a minute-by-minute census of Binance's monthly against daily archives.

\section{Related literature}

A large literature asks whether return predictability survives out of sample. \citet{welchgoyal2008} show that most equity-premium predictors fail out of sample, and \citet{campbellthompson2008} propose the out-of-sample $R^2$ and argue that small values can matter economically. At intraday horizons, \citet{gao2018} document intraday momentum in equity index returns, and \citet{chordia2008} show that short-horizon predictability from order imbalance declines as liquidity improves. For cryptocurrencies, \citet{liutsyvinski2021} find that returns are predicted by cryptocurrency-specific factors such as momentum and investor attention, and \citet{makarovschoar2020} document large price deviations across exchanges. We use only public candles from one venue, work at the 15-minute horizon, and focus on the gap between statistical and economic significance under a design that rules out specification search on the evaluation data.

On methods, we use the nested-model comparison of \citet{clarkwest2007} and \citet{clarkmccracken2001}, the equal-accuracy test of \citet{dieboldmariano1995}, Bartlett-kernel variances \citep{neweywest1987} with the bandwidth rule and the equal-weighted cosine estimator of \citet{lazarus2018}, the step-down procedure of \citet{holm1979}, block resampling \citep{kunsch1989}, and the stability tests of \citet{giacominiwhite2006} and \citet{giacominirossi2010}. On research practice, we follow the call of \citet{harvey2017} for pre-specification and for reporting null results.

\section{Data}

\subsection{Source and sample}
All market data are Binance Spot BTCUSDT one-minute candles from the exchange's public archive. Each candle contains open, high, low and close prices, base and quote volume, the number of trades and taker-buy volume. Table~\ref{tab:sample} lists the periods. The development sample was used to choose the primary test, the secondary specification and the power planning, and it has been analyzed repeatedly; any statement about it is exploratory. The historical archives were downloaded only after the historical protocol was frozen. The first holdout was opened only for the confirmatory test. A second holdout remains sealed, and the 1 September 2026 file, which served as a one-day engineering smoke test, is excluded from all statistics.

\begin{table}[ht]\centering\small
\caption{Sample periods and roles}\label{tab:sample}
\begin{tabular}{@{}llrl@{}}
\toprule
Period (UTC) & Role & Minutes & Status at design freeze \\
\midrule
2020-01 -- 2020-12 & Warm-up / training only & 527,040$^{a}$ & not downloaded \\
2021-01 -- 2025-08 & Historical evaluation (56 folds) & 2,453,760$^{a}$ & not downloaded \\
2025-09 -- 2026-06 & Development & 436,320 & seen (exploratory) \\
2026-07 -- 2026-08 & Holdout A (confirmatory) & 89,280 & sealed \\
2026-09-01 & Engineering smoke day & 1,440 & not used for statistics \\
2026-09-02 -- 2026-10-31 & Holdout B (confirmatory) & 86,400 & sealed \\
\bottomrule
\multicolumn{4}{@{}p{0.92\linewidth}@{}}{\footnotesize $^{a}$Calendar minutes (2,980,800 in total). The historical archives have 15 gaps that remove 2,325 minutes, leaving 2,978,475. No gaps were filled, no rows were dropped, and no feature or label spans a gap.}
\end{tabular}
\end{table}

\subsection{A census of Binance monthly and daily archives}\label{sec:census}
Before running any forecast we checked the historical data minute by minute. Binance publishes the same candles in monthly and in daily archives. A public report (issue \#475 of the \texttt{binance-public-data} repository) noted that some monthly files disagree with the corresponding daily files, and the December 2020 monthly file used here is byte-identical to the file described there. We therefore downloaded the daily file for every one of the 2,070 days of the historical sample and compared all fields of every minute. Every file matched its official SHA-256 checksum. Table~\ref{tab:census} reports the result.

\begin{table}[ht]\centering\small
\caption{Census of Binance monthly against daily one-minute archives, 2020-01-01 to 2025-08-31}\label{tab:census}
\begin{tabular}{@{}lrrr@{}}
\toprule
 & Selected days & Other days & All days \\
\midrule
Days verified & 11 & 2,059 & 2,070 \\
Minutes with field conflicts & 112 & 23 & 135 \\
Monthly-only minutes (zero-trade placeholders) & 22 & 0 & 22 \\
Total conflicting minutes & 134 & 23 & 157 \\
Minutes missing in both files & 569 & 1,756 & 2,325 \\
\bottomrule
\multicolumn{4}{@{}p{0.92\linewidth}@{}}{\footnotesize The 11 selected days were chosen because external reports or our result-blind checks flagged them, so they cannot be used to estimate a conflict rate. Whenever both files contain a conflicting minute, the daily file never shows fewer trades or less volume than the monthly file. In the 2,059 other days the 23 conflicting minutes form 21 episodes of at most two consecutive minutes; the largest price difference is 5.5 basis points.}
\end{tabular}
\end{table}

The conflicts are rare but not trivial. On 11 January 2021, eight minutes appear in the monthly file as zero-trade bars with equal open, high, low and close, while the daily file records 2,967 to 7,820 trades per minute and prices up to 332 basis points away. On 22--23 April 2021, monthly volume is up to about 90 times lower than daily volume. The monthly files also contain 211 zero-trade, zero-volume bars in four blocks. Two blocks appear identically in the daily files; the other two are among the conflicts above. We cannot tell which archive is correct without a third source. Under the frozen data dependencies of the program (a 60-minute feature window and a 15-minute label window), 1,457 of the 2,452,148 structurally eligible evaluation samples (0.059\%; counted before month-end label-maturity exclusions) touch a conflicting minute; a looser upper bound gives 0.070\%. The primary analysis uses the monthly files exactly as downloaded, as frozen in the protocol. Section~\ref{sec:sdq} reports a pre-registered sensitivity analysis that removes every affected sample.

\section{The frozen forecasting program}

\subsection{Timing, label and features}
Time stamps refer to the opening minute of each candle. A candle that opens at $t$ is complete at $t+60$s; its features become available at $t+61$s, the decision time. The label at horizon $h\in\{5,15\}$ minutes is the log close-to-close return $y_{t,h}=\log C_{t+h}-\log C_t$. It is observable one second after the closing time of the candle that opens at $t+h$, and it enters a training set only if it is observable at the fitting time. Features and labels are computed once on the concatenated series and never across an archive gap. Table~\ref{tab:features} lists the fourteen features (specification v1). Specification v2 drops \texttt{log\_close\_open}, which in the development data is almost collinear with the one-minute return ($\rho=0.99993$). Because that choice used development data, v2 is a secondary specification only.

\begin{table}[ht]\centering\small
\caption{Features (specification v1), computed at the close of candle $t$}\label{tab:features}
\begin{tabular}{@{}ll@{}}
\toprule
Feature & Definition \\
\midrule
\texttt{log\_return\_$k$m}, $k\in\{1,5,15,60\}$ & $\log C_t-\log C_{t-k}$ \\
\texttt{return\_std\_$w$m}, $w\in\{5,15,60\}$ & sample s.d.\ of the last $w$ one-minute log returns \\
\texttt{volume\_change\_1m} & $V_t/V_{t-1}-1$ \\
\texttt{volume\_zscore\_60m} & $(V_t-\bar V_{t-60:t-1})/s(V_{t-60:t-1})$ \\
\texttt{taker\_buy\_volume\_ratio} & taker-buy base volume$_t/V_t$ \\
\texttt{trades\_change\_60m} & $N_t/\bar N_{t-60:t-1}-1$ \\
\texttt{log\_high\_low} & $\log H_t-\log L_t$ \\
\texttt{log\_close\_open} & $\log C_t-\log O_t$ \\
\texttt{close\_range\_position} & $(C_t-L_t)/(H_t-L_t)$ \\
\bottomrule
\multicolumn{2}{@{}p{0.8\linewidth}@{}}{\footnotesize $C,O,H,L$: close, open, high, low; $V$: base volume; $N$: number of trades. Undefined ratios (zero volume, zero range) make the sample invalid; they are never filled.}
\end{tabular}
\end{table}

\subsection{Model and training scheme}
Each fit standardizes the features on the legal training rows only and estimates a Ridge regression \citep{hoerlkennard1970} with penalty $\alpha=1$ and an intercept. Nothing is tuned or selected at any point. Training windows expand from the start of each sample. In the historical test there is one fit per evaluation month, at 00:00 UTC on its first day, using between 0.53 and 2.93 million training samples, and the forecasts are evaluated over that month. The Holdout A test follows its frozen protocol: fits on 1 July and 1 August 2026, trained on all data from September 2025 onward. Holdout B fits on 2 September and 1 October 2026 with the same training start, so its training data include Holdout A. The historical and holdout tests therefore use different training histories.

\section{Inference}

\subsection{Primary test}
Let $f_t$ be the Ridge forecast and $b_t$ a benchmark. The \citet{clarkwest2007} adjusted differential is
\begin{equation}
c_t=(y_t-b_t)^2-\left[(y_t-f_t)^2-(f_t-b_t)^2\right]=2(f_t-b_t)(y_t-b_t),
\end{equation}
and the null of no predictive information is $H_0:\mathbb{E}[c_t]\le0$ against $H_1:\mathbb{E}[c_t]>0$. With $b_t=0$, $c_t=2f_ty_t$. Within each month we compute the mean of $c_t$ and a Bartlett long-run variance with lag $L_m=\max\{60,\lceil1.3\sqrt{T_m}\rceil\}$, the bandwidth rule of \citet{lazarus2018}, where $T_m$ is the number of valid evaluation minutes. When archive gaps or invalid samples split a month, the variance is computed within each contiguous segment and the segments are combined with zero cross-segment covariance. This is a modelling assumption, fixed before any forecast existed. Months are pooled with weights proportional to $T_m$ and zero cross-month covariance, and the $p$-value uses normal critical values. The null is rejected if $p<0.05$. The primary family contains only this test: specification v1, $h=15$, benchmark zero.

\subsection{Secondary tests and pre-specified robustness}
Five secondary tests form a \citet{holm1979} family: (S1) specification v2 against zero; (S2) the 5-minute horizon; (S3) the historical-mean benchmark; (S4) the Diebold--Mariano squared-error test, two-sided; and (S5) the time-series information coefficient, a HAC test on monthly centred and standardized products of forecasts and outcomes. Secondary results cannot upgrade a non-rejection of the primary test. For the historical test, three robustness checks for cross-month dependence were specified in advance: a HAC test on the series of 56 monthly means (lag 4); a circular block bootstrap of months with block lengths 3, 6 and 12 \citep{kunsch1989}; and a comparison of equal with sample-size weighting. Two placebos were also pre-specified as exploratory diagnostics. P1 flips forecast signs in random 60-minute blocks (500 draws). P2 pairs each forecast with the outcome exactly one day later.

\subsection{Data-quality sensitivity}\label{sec:sdq}
Before any real forecast was produced, we fixed a sensitivity analysis that removes from training and evaluation every sample whose feature or label window contains any of the 157 conflicting minutes, and then re-estimates the model. A short-segment rule, also fixed in advance, drops any evaluation segment that the mask leaves no longer than the HAC lag (one 15-minute segment in December 2021). If this analysis and the primary test disagree, the conclusion is labelled data-source-sensitive.

\subsection{Exploratory supplementary analyses}
After the historical and Holdout A tests were complete, and before Holdout B was run, we ran further analyses requested in an internal review. They use only the stored forecasts and outcomes, fit no model, and are labelled exploratory. They cover drift and timing, alternative variance estimators and subsamples, the economic magnitude of the forecasts, and confidence bounds and calibration. Our implementation reproduces the frozen primary statistics exactly, and two independent runs produced byte-identical outputs (\nolinkurl{docs/supplementary-evidence/determinism-check.json}).

\section{Research protocol}

Table~\ref{tab:timeline} lists the freezes. Each frozen document and each source file of the forecasting program is identified by its SHA-256 hash, and the hashes were checked before every real-data process. The number of real-data processes was capped and each was logged; a crashed formal run could not be repeated without a new decision. The historical and holdout tests were each completed twice in independent processes, and all deterministic outputs were byte-identical (26, 9 and 9 files for the historical test and the two holdouts; for Holdout B after one crashed, result-blind formal run, Section~\ref{sec:holdoutb}). Before any fitting, a result-blind audit reconciled every sample count against an independent minute-grid calculation. The archive census was a separate data-quality exercise. Its downloads were made once and its execution deviated from its own protocol in logged ways. Its comparison and aggregation were replayed with byte-identical output (17 files), its main part was recomputed independently, and the final 14-day supplement received only a light audit.

\begin{table}[ht]\centering\small
\caption{Timeline of freezes and data access (2026)}\label{tab:timeline}
\begin{tabular}{@{}lp{0.70\linewidth}@{}}
\toprule
Date & Event \\
\midrule
2 Nov & Forecasting program fixed (features, label, timing, model) \\
4 Nov & Holdout A protocol frozen and hashed; code guards seal both holdouts \\
5 Nov & Historical-extension protocol frozen before any historical download \\
6--9 Nov & Historical archives downloaded and audited (no forecasts) \\
10--14 Nov & Archive census against daily files (Section~\ref{sec:census}) \\
14 Nov & Historical execution contract and data-quality sensitivity frozen \\
16 Nov & Historical test executed and reproduced \\
16--17 Nov & Holdout A contract frozen; written authorization citing the protocol hash; download, audit, execution, reproduction \\
17 Nov & Exploratory supplementary analyses; Holdout B protocol frozen; Holdout B authorized, downloaded and audited; first formal run stopped (result-blind); re-audit, execution and reproduction under a recorded owner decision (Section~\ref{sec:holdoutb}) \\
\bottomrule
\end{tabular}
\end{table}

Some exposures cannot be removed. The holdout months had already passed when the design was frozen: they were sealed, not in the future, and general knowledge of market conditions in 2026 cannot be excluded. An early metadata request recorded the lengths and checksums of the Holdout A archives, though not their contents. The engineering smoke day of 1 September 2026 fixes the price level at the end of Holdout A and the start of Holdout B. The historical results were known when the Holdout A test was run. Its primary and secondary tests had been fixed earlier, but its execution contract, which added procedural rules and a descriptive drift decomposition, was frozen after the historical results were known. The Holdout A archives were not cross-checked against daily files, but they matched byte for byte the checksums and file lengths recorded two weeks earlier. The Holdout B authorization was given in reply to a request that quoted the protocol hash, rather than by repeating the hash, and its first formal run stopped before writing any result; it was repeated once under a recorded owner decision (Section~\ref{sec:holdoutb}). The Holdout B protocol, including the choice of the information coefficient as its primary statistic, was written after the historical, Holdout A and exploratory supplementary results were known. The Holdout B archives were not cross-checked against daily files either; they were verified only against Binance's official CHECKSUM files.

\section{Results}

\subsection{Historical external validation, 2021--2025}
Table~\ref{tab:hist} reports the pre-registered results. The primary test rejects the null of no predictive information: the mean differential is $3.02\times10^{-8}$ with a standard error of $5.54\times10^{-9}$ ($t=5.44$, $p=2.6\times10^{-8}$). Four of the five secondary tests reject after the Holm adjustment. The exception is the Diebold--Mariano test: the mean squared error falls by $8.1\times10^{-9}$, an out-of-sample $R^2$ of $0.071\%$, which is not significant ($p=0.125$). None of the pre-specified dependence flags is triggered. The test on the 56 monthly means gives $t=2.14$ ($p=0.016$); the 95\% block-bootstrap intervals exclude zero for block lengths 3, 6 and 12; and equal and sample-size weighting give $t=5.49$ and $5.44$. The data-quality sensitivity gives $t=5.36$, so the conclusion does not depend on the archive conflicts.

\begin{table}[ht]\centering\small
\caption{Historical external validation: 56 monthly folds, 2021-01 to 2025-08 (pre-registered)}\label{tab:hist}
\begin{tabular}{@{}lrrrrl@{}}
\toprule
Test & Mean & $t$ & $p$ & Holm $p$ & Decision \\
\midrule
Primary: v1, 15m, CW vs.\ zero & $3.02\times10^{-8}$ & 5.44 & $2.6\times10^{-8}$ & -- & reject \\
S1: v2, 15m, CW vs.\ zero & $3.03\times10^{-8}$ & 5.43 & $2.8\times10^{-8}$ & $5.6\times10^{-8}$ & reject \\
S2: v1, 5m, CW vs.\ zero & $1.28\times10^{-8}$ & 9.31 & $6.6\times10^{-21}$ & $3.3\times10^{-20}$ & reject \\
S3: v1, 15m, CW vs.\ historical mean & $3.01\times10^{-8}$ & 5.52 & $1.7\times10^{-8}$ & $5.0\times10^{-8}$ & reject \\
S4: v1, 15m, DM vs.\ zero (two-sided) & $-8.1\times10^{-9}$ & $-1.53$ & 0.125 & 0.125 & no rejection \\
S5: v1, 15m, information coefficient & 0.024 & 6.47 & $4.8\times10^{-11}$ & $1.9\times10^{-10}$ & reject \\
Data-quality sensitivity (S-DQ) & $2.98\times10^{-8}$ & 5.36 & $4.2\times10^{-8}$ & -- & reject \\
\midrule
\multicolumn{6}{@{}l}{Evaluation minutes: 2,451,190. Zero-forecast RMSE: 33.7 bp. Out-of-sample $R^2$ vs.\ zero: $0.071\%$.}\\
\bottomrule
\multicolumn{6}{@{}p{0.95\linewidth}@{}}{\footnotesize Mean: pooled mean of the Clark--West differential (primary, S1--S3, S-DQ), of $(y-f)^2-(y-b)^2$ (S4; negative values favour the model) or of the information coefficient (S5). S1 uses a specification chosen with development data.}
\end{tabular}
\end{table}

\paragraph{How strong is the evidence?} Table~\ref{tab:sens} collects the exploratory sensitivity results. Alternative variance estimators move the $t$-statistic between 4.28 (Bartlett lag 60) and 5.86 (the equal-weighted cosine estimator with fixed-$b$ Student-$t$ critical values, which drops the zero cross-month covariance assumption). The more conservative test on monthly means gives 2.14, which indicates that the monthly estimates vary more across months than the within-month variance implies. The evidence is concentrated in time: 2021 contributes 76\% of the summed differential, and the yearly $t$-statistics are 4.59, 1.05, 2.05, 2.42 and 1.29 for 2021--2025. Excluding 2021, the statistic is 3.00 ($p=0.0013$) and the monthly-means test gives 4.82. A Giacomini--Rossi fluctuation test on 12-month rolling windows rejects stability: the maximum statistic, 3.53, occurs in the window ending December 2021, against a simulated 5\% critical value of 3.22. The rolling mean is nevertheless positive in every window. The result is not driven by outliers. Removing the 10, 100 or 1,000 most extreme minutes, winsorizing, or excluding five stress windows leaves $t$ between 5.8 and 6.4. A conditional test of equal squared-error accuracy \citep{giacominiwhite2006}, with the previous month's realized volatility as instrument, finds no state dependence (joint Wald $p=0.44$).

\begin{table}[ht]\centering\small
\caption{Robustness of the primary statistic (exploratory)}\label{tab:sens}
\begin{tabular}{@{}lrr@{}}
\toprule
 & History & Holdout A \\
\midrule
Frozen: Bartlett within month, $L=\max(60,\lceil1.3\sqrt{T_m}\rceil)$ & 5.44 & 1.89 \\
Bartlett within month, $L=60$ & 4.28 & 2.01 \\
Bartlett within month, twice the frozen lag & 5.67 & 1.87 \\
Bartlett on the concatenated series (no zero cross-month covariance) & 5.32 & 1.87 \\
Equal-weighted cosine, fixed-$b$ $t$ critical values (concatenated) & 5.86 & 1.93 \\
HAC on monthly means (pre-specified S\_dep1) & 2.14 & -- \\
Excluding 2021 / 2022--2025 only & 3.00 & -- \\
Within-month demeaned CW (timing only) & 5.74 & 3.48 \\
Real-time demeaned CW (past-week means, mature labels only) & 5.65 & 3.15 \\
CW against the past-week mean return & 9.07 & 2.96 \\
Volatility-standardized CW (past-week s.d.) & 2.76 & 0.14 \\
Placebo P1, block sign flips: empirical $p$ & 0.002 & 0.028 \\
Shift placebo: share of 174 / 114 shifts with $t\ge$ observed & 0 & 0 \\
\bottomrule
\multicolumn{3}{@{}p{0.9\linewidth}@{}}{\footnotesize Entries are $t$-statistics unless stated otherwise. Rows below the first are exploratory and were computed from the stored forecasts after the historical and Holdout A tests.}
\end{tabular}
\end{table}

\paragraph{Drift and timing.} With a zero benchmark, $c_t=2f_ty_t$ also rewards a forecast whose average agrees with the sample's average return, even when its timing is wrong. The pre-specified one-day shift placebo P2 was mildly significant ($p\approx0.03$--0.04), which raised this concern. Decomposing the mean differential into $2\bar f\bar y$ and $2\,\mathrm{cov}(f,y)$ within each month settles it. The drift term is negative in every year. The covariance (timing) term accounts for more than the entire statistic: $3.11\times10^{-8}$ against $3.02\times10^{-8}$. A within-month demeaned test, which removes drift by construction, gives $t=5.74$; excluding 2021 it gives 3.36. Real-time demeaning with past-week means gives 5.65. Over a distribution of 174 shifts of at least one day, the observed statistic exceeds every shifted value. Whole-day shifts average $t=0.22$, and 2 of 21 exceed 1.645, a little above the nominal 5\%, which suggests a weak time-of-day effect rather than drift.

\paragraph{What the model uses (descriptive).} The coefficients on the 5- and 60-minute past returns are negative in every historical fit, and the one on the 15-minute return in 55 of 56: the program is mainly a short-term reversal forecast. The coefficient on 15-minute realized volatility is positive in every month. The forecasts degrade with execution delay. Pairing the forecast with the outcome that starts one minute later lowers the statistic from 5.44 to 5.19, and five minutes later to 3.78. The signal is therefore not only a bounce between the last trade of one candle and the first trade of the next.

\subsection{Confirmatory holdout, July--August 2026}
Table~\ref{tab:holdout} reports the confirmatory test. Both months have 44,625 valid evaluation minutes, with no gaps, anomalous candles or invalid samples. The pooled statistic is $t=1.89$ (one-sided $p=0.029$), so the pre-registered primary test rejects the null. The evidence is much weaker than in the historical test: July alone is not significant ($p=0.20$), August is borderline ($p=0.043$), and the rejection holds under every variance estimator in Table~\ref{tab:sens} ($t$ between 1.87 and 2.01) and under the block sign-flip placebo ($p=0.028$). Among the secondary tests, the 5-minute horizon, the historical-mean benchmark and the information coefficient reject after the Holm adjustment; specification v2 does not.

\begin{table}[ht]\centering\small
\caption{Confirmatory Holdout A test under the frozen protocol (July--August 2026)}\label{tab:holdout}
\begin{tabular}{@{}lrrrr@{}}
\toprule
 & Mean & s.e. & $t$ & one-sided $p$ \\
\midrule
Primary, July 2026 & $2.78\times10^{-9}$ & $3.32\times10^{-9}$ & 0.84 & 0.201 \\
Primary, August 2026 & $8.10\times10^{-9}$ & $4.71\times10^{-9}$ & 1.72 & 0.043 \\
Primary, pooled & $5.44\times10^{-9}$ & $2.88\times10^{-9}$ & 1.89 & 0.029 \\
\midrule
 & Raw $p$ & Holm $p$ & \multicolumn{2}{l}{Decision} \\
\midrule
S1: v2, 15m, CW vs.\ zero & 0.039 & 0.079 & \multicolumn{2}{l}{no rejection} \\
S2: v1, 5m, CW vs.\ zero & 0.0033 & 0.013 & \multicolumn{2}{l}{reject} \\
S3: v1, 15m, CW vs.\ historical mean & 0.0061 & 0.018 & \multicolumn{2}{l}{reject} \\
S4: v1, 15m, DM vs.\ zero (two-sided) & 0.955 & 0.955 & \multicolumn{2}{l}{no rejection} \\
S5: v1, 15m, information coefficient (0.035) & 0.00046 & 0.0023 & \multicolumn{2}{l}{reject} \\
\midrule
\multicolumn{5}{@{}l}{Evaluation minutes: 89,250. Zero-forecast RMSE: 18.4 bp. Out-of-sample $R^2$ vs.\ zero: $-0.005\%$.}\\
\bottomrule
\end{tabular}
\end{table}

The holdout forecasts have a negative average ($-0.35$ basis points), partly inherited from the intercept (about $-0.2$ basis points) of a training period with falling prices. Realized returns had a positive average over the two months, so the drift term works \emph{against} the forecasts. The drift-free tests are correspondingly stronger than the primary test: the within-month demeaned statistic is $t=3.48$ ($p=0.0003$), and real-time demeaning gives 3.15. The holdout signal is concentrated in volatile minutes: standardizing by past-week volatility removes it ($t=0.14$). It also decays quickly with execution delay ($t=0.95$ after one minute). The holdout model differs from the historical models. Because its training data begin in September 2025 (the development period, plus July 2026 for the August fit), its two largest coefficients belong to the near-duplicate pair of features, \texttt{log\_return\_1m} and \texttt{log\_close\_open}. The coefficients are almost equal and of opposite sign, so their net effect loads on the gap between the previous close and the current open. A training-period audit traced this leverage to a single crash session.

\subsection{How large is the information?}
\paragraph{Accuracy bounds.} In the historical sample and in Holdout A the 95\% one-sided upper bound on the out-of-sample $R^2$ is about $0.15\%$, and an equivalence test rejects $R^2\ge0.5\%$ ($p<10^{-7}$). Squaring the information coefficient gives the $R^2$ that an optimally scaled forecast could reach: $0.055\%$ historically (95\% interval $0.027$--$0.094\%$) and $0.12\%$ in the holdout ($0.02$--$0.31\%$). Mincer--Zarnowitz regressions explain why the Clark--West test rejects while the squared-error test does not. Historically the forecast slope is 0.68 (s.e.\ 0.11): the forecasts are somewhat too dispersed. In the holdout the slope is 1.01 (0.33) but the intercept is 0.84 basis points (0.26). This bias costs more squared error than the signal can save. Rescaling the forecasts in real time, with a slope estimated only from earlier months, raises the out-of-sample $R^2$ over February 2021--August 2025 (January 2021 has no earlier month) from $0.057\%$ to $0.067\%$.

\paragraph{Gross economics.} Table~\ref{tab:econ} summarizes the forecasts' economic scale. We do not backtest a strategy: there is no model of latency, queue position, adverse selection or capacity. The typical absolute 15-minute forecast is 0.58 basis points historically and 0.44 in the holdout. Only 7.9\% (history) and 1.3\% (holdout) of minutes have forecasts above 2 basis points. Minutes with larger forecasts do have larger forecast-signed returns, with hit rates rising from 52\% to 55\%. We evaluate a non-overlapping rule that, every 15 minutes, takes the sign of the forecast when its absolute value exceeds $k$ basis points. The break-even fee per side is 0.21 basis points when the rule always trades. It rises to 0.85 basis points for $k=2$ and to about 4.6 basis points for $k=5$, where the rule is active in only 1.3\% of periods and the confidence interval is wide. Binance's standard (VIP 0) spot taker fee is 10 basis points, or 7.5 with its fee-token discount. The quoted spread of BTCUSDT is typically one price tick of 0.01 USDT, well below 0.01 basis points at 2021--2026 price levels (an exchange specification, not measured here), so fees, latency and adverse selection are the binding costs. Only very high-volume fee tiers approach the break-even levels of the most selective rule, and even then latency and adverse selection would remain. A reader who maps the historical $R^2$ into a hypothetical continuously rebalanced portfolio, as in \citet{campbellthompson2008}, would obtain a gross annualized Sharpe ratio of about 5. That figure assumes trading every minute with no costs. It comes almost entirely from 2021, and the implementable non-overlapping rule has a gross Sharpe ratio of 0.8.

\begin{table}[ht]\centering\footnotesize
\caption{Economic scale of the 15-minute forecasts (exploratory, gross of all costs)}\label{tab:econ}
\begin{tabular}{@{}lrrr@{}}
\toprule
 & History & History 2022--25 & Holdout A \\
\midrule
Median $|f|$ (bp) / 99th percentile (bp) & 0.58 / 5.5 & -- & 0.44 / 2.2 \\
Forecast-signed return, all minutes (bp) & 0.21 & -- & 0.14 \\
Forecast-signed return, $|f|>2$ bp (bp) / hit rate & 1.45 / 54.4\% & -- & 3.01 / 54.7\% \\
\midrule
\multicolumn{4}{@{}l}{Break-even fee per side (bp) of the non-overlapping rule, with 95\% month-bootstrap interval:}\\
$k=0$ (always trade) & 0.21 [0.03, 0.40] & 0.15 [$-$0.07, 0.39] & 0.28 [$-$0.22, 0.79] \\
$k=2$ bp (trades 8 / 5 / 2\%) & 0.85 [0.13, 1.53] & 0.34 [$-$0.53, 1.22] & 1.25 [0.85, 1.46] \\
$k=5$ bp (trades 1.3 / 0.6 / 0.15\%) & 4.6 [1.5, 7.4] & 3.4 [0.4, 6.5] & 5.6 [$-$1.1, 7.8] \\
\midrule
Reference: Binance VIP 0 spot taker fee & \multicolumn{3}{c}{10 bp (7.5 bp with fee-token discount)} \\
\bottomrule
\multicolumn{4}{@{}p{0.9\linewidth}@{}}{\footnotesize The rule decides every 15 minutes at a fixed phase; across the 15 possible phases the always-trade break-even ranges from $-0.28$ to 0.97 bp in the history and from $-0.19$ to 0.72 bp in Holdout A. For the 5-minute horizon the always-trade break-even is 0.15 bp historically and 0.07 bp in Holdout A, where the rule has no positive break-even for $k\ge1$.}
\end{tabular}
\end{table}

\section{Second confirmatory holdout: Holdout B}\label{sec:holdoutb}
Holdout B covers 2 September to 31 October 2026. Its protocol was frozen before any of its data were accessed. The file \nolinkurl{docs/confirmatory-protocol-holdout-b.md} has SHA-256
\begin{center}\footnotesize\texttt{6843156fdcd455956a3dc9cbc27cb5c2}\\\texttt{13cf22b2fa7932ca49e6c51e13e2f315}\end{center}
Because the 1 September smoke file fixes the starting price level, the primary statistic is drift-free: the within-month time-series information coefficient of the v1 15-minute forecasts, tested one-sided at 5\% with the frozen variance and pooling rules. Training expands from September 2025 and includes Holdout A, and the rows of 1 September 2026 were discarded at parsing. Ex ante, the test had power of about 0.60 if the true effect equals the 2022--2025 historical value and about 0.95 at the Holdout A value.

The first formal run stopped in the pooling step before writing any result; only a stop file and a traceback were recorded. The runner labelled each fold by the month of its first evaluation bar, and October's first evaluation bar opens at 23:59 on 30 September. No statistic had been written or displayed. Under a recorded owner decision, the label was changed to the fit-time month, as in the Holdout A runner. The data audit was repeated and matched the original in every field except the recorded hash of the execution code. The formal run and its reproduction were then run once each, with byte-identical outputs.

\begin{table}[ht]\centering\small
\caption{Second confirmatory holdout (Holdout B, September--October 2026)}\label{tab:holdoutb}
\begin{tabular}{@{}lrrrr@{}}
\toprule
Primary: v1, 15m, information coefficient & IC & s.e. & $t$ & one-sided $p$ \\
\midrule
September 2026 (from 2 September) & 0.025 & 0.026 & 0.95 & 0.171 \\
October 2026 & 0.018 & 0.020 & 0.93 & 0.175 \\
Pooled & 0.021 & 0.016 & 1.33 & 0.092 \\
\midrule
 & Raw $p$ & Holm $p$ & \multicolumn{2}{l}{Decision} \\
\midrule
S1: v1, 15m, CW vs.\ zero & 0.093 & 0.371 & \multicolumn{2}{l}{no rejection} \\
S2: v1, 15m, within-month demeaned CW & 0.093 & 0.371 & \multicolumn{2}{l}{no rejection} \\
S3: v1, 15m, CW vs.\ historical mean & 0.098 & 0.371 & \multicolumn{2}{l}{no rejection} \\
S4: v1, 15m, DM vs.\ zero (two-sided) & 0.914 & 0.914 & \multicolumn{2}{l}{no rejection} \\
S5: v1, 5m, information coefficient & 0.032 & 0.159 & \multicolumn{2}{l}{no rejection} \\
\midrule
\multicolumn{5}{@{}l}{Evaluation minutes: 86,309. Out-of-sample $R^2$ vs.\ zero: $0.012\%$ (95\% interval $-0.21$ to $0.23\%$).}\\
\bottomrule
\end{tabular}
\end{table}

The primary test does not reject: the pooled information coefficient is 0.021 (s.e.\ 0.016, $t=1.33$, $p=0.092$). No secondary test rejects after the Holm adjustment; the smallest raw $p$ is 0.032, for the 5-minute horizon. The drift term is small and of opposite sign in the two months, so it nearly cancels in the pooled statistic (the raw and demeaned Clark--West tests give the same $p$), and the always-trade rule breaks even at 0.29 basis points per side. The estimate is close to the historical values (0.0235 for 2021--2025 and 0.0198 for 2022--2025). A real but small effect would often produce a result like this in two months, but the result is equally compatible with no effect. As pre-specified, it is reported as a failed confirmatory replication. Pooling the two holdouts gives an information coefficient of 0.028 ($t=2.95$); this pooling is descriptive and is not a test.

\section{Discussion}

The evidence on predictive information is strong over the historical sample, holds at the pre-registered level in the first holdout and is not confirmed in the second. Two months are too few to detect an information coefficient of about 0.02 reliably: the second holdout's ex-ante power was about 0.6 at that size (computed with the Holdout A standard error; the realized standard error, 0.016, was larger, so realized power was lower), and its estimate is compatible both with the historical value and with zero. Historically it holds across variance estimators and drift-free variants and is robust to source defects and outliers; in the first holdout it holds across variance estimators and drift-free variants, while in the second none of the drift-free statistics (the information coefficient, the demeaned Clark--West statistic and the v2 information coefficient) reaches the 5\% level. Two qualifications matter. First, the strength of the historical evidence depends on how cross-month variation is treated: within-month HAC gives $t=5.44$, while the monthly-means test gives 2.14. Second, the historical evidence is concentrated in 2021, a year of exceptional volatility. In economic terms the evidence is a negative result. The forecasts reduce squared error by at most a fraction of a percent, and their gross edge is below the fees faced by ordinary traders even for the most selective rule (a break-even of about 3.4 to 5.6 basis points per side in the historical sample and Holdout A, and negative in Holdout B, against 7.5 to 10), before latency and adverse selection. A predictable component of this size is consistent with a market that is efficient at the margin that matters \citep{chordia2008}.

The contrast between the evidence layers carries a methodological lesson. The Clark--West test is the right instrument for asking whether a nested model carries information, and with weak signals it is far more powerful than squared-error comparisons \citep{clarkwest2007,clarkmccracken2001}. But a rejection is not evidence of lower forecast error. With a zero benchmark it can also reward agreement with the sample drift, and a forecast bias can erase the accuracy gain. Reporting drift-free variants, time-shift placebos, calibration regressions and gross economics next to the primary statistic makes clear which kind of information drives a rejection.

\paragraph{Limitations.} Each holdout covers only two months and has modest power, and the two holdouts point in different directions. The historical test evaluates a design created in 2026 on earlier data. It is an external validation, not a prospective test, and general knowledge of the period cannot be excluded. The results are for one instrument (BTCUSDT spot) on one venue, with one model class and two horizons. The historical and holdout tests use different training histories (the holdouts train from September 2025 only, and Holdout B's training includes Holdout A). The within-month variance assumes independence across archive gaps. The economic analysis is gross of all costs and ignores execution. The census shows that the two Binance archives disagree, but it does not show which one is right. The supplementary analyses are exploratory and were chosen after the confirmatory results were known.

\section{Conclusion}

A simple, fully pre-specified program that uses only public one-minute candles of BTCUSDT carries out-of-sample information about 15-minute returns. The evidence is strong but uneven over 56 historical months. It holds, narrowly, in the first confirmatory holdout and is not confirmed in the second, where the estimate is similar in size but not significant. The information concerns timing, not drift. It is economically small: it yields no reliable accuracy gain, and its gross edge lies below ordinary trading costs. Frozen protocols, result-blind data audits, byte-identical reproduction and a full account of what was known before each test are practical with public high-frequency data. With them, a statistically significant result can be reported together with its economic size.

\section*{Disclosure statement}
The author has no financial relationship with Binance or any trading firm and reports no conflicts of interest. The study involves no trading, no account access and no non-public data. \emph{Use of AI tools.} The author used AI assistants (Anthropic's Claude and OpenAI's Codex), under the author's direction, to write and test code, run the pre-specified computations and supplementary analyses, check results against the evidence files, and draft and edit the text. All research questions, design choices, data-access and holdout decisions, and interpretations were made or approved by the author, who takes full responsibility for the content.

\section*{Data, code and license}
Raw market data are not redistributed, in line with the Binance Vision dataset terms. The repository contains the frozen protocols, execution contracts, the hashes of every input file, download ledgers, the archive census, all result files and the scripts needed to rebuild every number in this paper from the public archives. Software is licensed under the PolyForm Noncommercial License 1.0.0. The paper and documentation are licensed under CC BY-NC-ND 4.0, and evidence derived from Binance data under CC BY-NC-SA 4.0. Commercial use requires a separate written license. Any use must cite this paper.

\begin{thebibliography}{99}
\setlength{\itemsep}{1pt}\small
\bibitem[Campbell and Thompson(2008)]{campbellthompson2008} Campbell, J.~Y., and S.~B. Thompson (2008). Predicting excess stock returns out of sample: Can anything beat the historical average? \emph{Review of Financial Studies} 21(4), 1509--1531.
\bibitem[Chordia et al.(2008)]{chordia2008} Chordia, T., R. Roll, and A. Subrahmanyam (2008). Liquidity and market efficiency. \emph{Journal of Financial Economics} 87(2), 249--268.
\bibitem[Clark and McCracken(2001)]{clarkmccracken2001} Clark, T.~E., and M.~W. McCracken (2001). Tests of equal forecast accuracy and encompassing for nested models. \emph{Journal of Econometrics} 105(1), 85--110.
\bibitem[Clark and West(2007)]{clarkwest2007} Clark, T.~E., and K.~D. West (2007). Approximately normal tests for equal predictive accuracy in nested models. \emph{Journal of Econometrics} 138(1), 291--311.
\bibitem[Diebold and Mariano(1995)]{dieboldmariano1995} Diebold, F.~X., and R.~S. Mariano (1995). Comparing predictive accuracy. \emph{Journal of Business \& Economic Statistics} 13(3), 253--263.
\bibitem[Gao et al.(2018)]{gao2018} Gao, L., Y. Han, S.~Z. Li, and G. Zhou (2018). Market intraday momentum. \emph{Journal of Financial Economics} 129(2), 394--414.
\bibitem[Giacomini and Rossi(2010)]{giacominirossi2010} Giacomini, R., and B. Rossi (2010). Forecast comparisons in unstable environments. \emph{Journal of Applied Econometrics} 25(4), 595--620.
\bibitem[Giacomini and White(2006)]{giacominiwhite2006} Giacomini, R., and H. White (2006). Tests of conditional predictive ability. \emph{Econometrica} 74(6), 1545--1578.
\bibitem[Harvey(2017)]{harvey2017} Harvey, C.~R. (2017). Presidential address: The scientific outlook in financial economics. \emph{Journal of Finance} 72(4), 1399--1440.
\bibitem[Harvey et al.(2016)]{harvey2016} Harvey, C.~R., Y. Liu, and H. Zhu (2016). \ldots\ and the cross-section of expected returns. \emph{Review of Financial Studies} 29(1), 5--68.
\bibitem[Hoerl and Kennard(1970)]{hoerlkennard1970} Hoerl, A.~E., and R.~W. Kennard (1970). Ridge regression: Biased estimation for nonorthogonal problems. \emph{Technometrics} 12(1), 55--67.
\bibitem[Holm(1979)]{holm1979} Holm, S. (1979). A simple sequentially rejective multiple test procedure. \emph{Scandinavian Journal of Statistics} 6(2), 65--70.
\bibitem[K\"unsch(1989)]{kunsch1989} K\"unsch, H.~R. (1989). The jackknife and the bootstrap for general stationary observations. \emph{Annals of Statistics} 17(3), 1217--1241.
\bibitem[Lazarus et al.(2018)]{lazarus2018} Lazarus, E., D.~J. Lewis, J.~H. Stock, and M.~W. Watson (2018). HAR inference: Recommendations for practice. \emph{Journal of Business \& Economic Statistics} 36(4), 541--559.
\bibitem[Liu and Tsyvinski(2021)]{liutsyvinski2021} Liu, Y., and A. Tsyvinski (2021). Risks and returns of cryptocurrency. \emph{Review of Financial Studies} 34(6), 2689--2727.
\bibitem[Makarov and Schoar(2020)]{makarovschoar2020} Makarov, I., and A. Schoar (2020). Trading and arbitrage in cryptocurrency markets. \emph{Journal of Financial Economics} 135(2), 293--319.
\bibitem[Newey and West(1987)]{neweywest1987} Newey, W.~K., and K.~D. West (1987). A simple, positive semi-definite, heteroskedasticity and autocorrelation consistent covariance matrix. \emph{Econometrica} 55(3), 703--708.
\bibitem[Welch and Goyal(2008)]{welchgoyal2008} Welch, I., and A. Goyal (2008). A comprehensive look at the empirical performance of equity premium prediction. \emph{Review of Financial Studies} 21(4), 1455--1508.
\bibitem[White(2000)]{white2000} White, H. (2000). A reality check for data snooping. \emph{Econometrica} 68(5), 1097--1126.
\end{thebibliography}

\end{document}
